%-------------------------%
% Custom commands

% RME:
% Requires: #1 contains a bullet body string that can be rendered by LaTeX.
% Modifies: The active itemize environment by adding one bullet item.
% Effects: Typesets a compact resume bullet.
% Inputs: #1 = bullet text.
% Outputs: One bullet item in the generated resume PDF.
\newcommand{\resumeItem}[1]{
  \item\small{#1}
}

% RME:
% Requires: #1-#4 contain organization, location, role or degree, and date text.
% Modifies: The active resume subheading list by adding one two-line entry.
% Effects: Typesets a primary education or experience heading with balanced spacing.
% Inputs: #1 = left title; #2 = right location; #3 = left subtitle; #4 = right date.
% Outputs: One primary heading block in the generated resume PDF.
\newcommand{\resumeSubheading}[4]{
  \vspace{2pt}\item
    \begin{tabular*}{0.97\textwidth}[t]{l@{\extracolsep{\fill}}r}
      \textbf{#1} & #2 \\
      \textit{\small#3} & \textit{\small #4} \\
    \end{tabular*}\vspace{-2pt}
}

% RME:
% Requires: #1 and #2 contain role and date text.
% Modifies: The active resume subheading list by adding one secondary entry.
% Effects: Typesets an additional role under the previous organization with readable separation.
% Inputs: #1 = left role text; #2 = right date text.
% Outputs: One secondary heading block in the generated resume PDF.
\newcommand{\resumeSubSubheading}[2]{
    \vspace{4pt}\item
    \begin{tabular*}{0.97\textwidth}{l@{\extracolsep{\fill}}r}
      \textit{\small#1} & \textit{\small #2} \\
    \end{tabular*}\vspace{-2pt}
}

% RME:
% Requires: #1 and #2 contain project title/details and optional right-side text.
% Modifies: The active resume subheading list by adding one project heading.
% Effects: Typesets a project heading with modest separation before its bullets.
% Inputs: #1 = project heading text; #2 = optional right-aligned text.
% Outputs: One project heading block in the generated resume PDF.
\newcommand{\resumeProjectHeading}[2]{
    \vspace{4pt}\item
    \begin{tabular*}{0.97\textwidth}{l@{\extracolsep{\fill}}r}
      \small#1 & #2 \\
    \end{tabular*}\vspace{-2pt}
}

% RME:
% Requires: #1 contains a short non-bulleted detail line.
% Modifies: The active education list by adding one aligned detail row.
% Effects: Renders supporting education details without opening a bullet list.
% Inputs: #1 = detail text.
% Outputs: Typeset detail line in the generated resume PDF.
\newcommand{\resumeEducationDetail}[1]{
    \vspace{1pt}\item[]
    \small{#1}\vspace{-2pt}
}

% RME:
% Requires: #1 contains a nested bullet body string.
% Modifies: The active itemize environment by adding one compact nested item.
% Effects: Typesets a compact nested resume item.
% Inputs: #1 = nested bullet text.
% Outputs: One nested bullet item in the generated resume PDF.
\newcommand{\resumeSubItem}[1]{\resumeItem{#1}\vspace{-2pt}}

\renewcommand\labelitemii{$\vcenter{\hbox{\tiny$\bullet$}}$}

\newcommand{\resumeSubHeadingListStart}{\begin{itemize}[leftmargin=0.15in, label={}, topsep=0pt, partopsep=0pt, parsep=0pt, itemsep=0pt]}
\newcommand{\resumeSubHeadingListEnd}{\end{itemize}\vspace{-1pt}}
\newcommand{\resumeItemListStart}{\begin{itemize}[leftmargin=0.22in, topsep=3pt, partopsep=0pt, parsep=0pt, itemsep=2pt]}
\newcommand{\resumeItemListEnd}{\end{itemize}\vspace{-2pt}}

\definecolor{Black}{RGB}{0, 0, 0}
\newcommand{\seticon}[1]{\textcolor{Black}{\csname #1\endcsname}}
